% SGA 6, Exposición XII: fragmento español alineado con la fuente.
% Alcance: apertura, Advertencia y sección 1 completas; termina antes de la sección 2.
% Fuente francesa: sga6_fr_r1/src/fr.tex, líneas 14169--14218.
% SHA-256 de la fuente francesa completa:
% 0D2158B9FE3B466309A92165727A0B191422E05FBAA0DACD5806E8322C54BA1A.
% Las superficies matemáticas de la fuente se conservan diplomáticamente;
% las anomalías detectadas se notifican por separado para su resolución canónica.

\section*{Exposición XII. Un teorema de representabilidad relativa sobre el funtor de Picard}
\begin{center}
por M. Raynaud (redactado por S. Kleiman)
\end{center}

\subsection*{Advertencia}

El objeto de las dos Exposiciones XII y XIII es dar demostraciones de los
nueve teoremas de finitud para el funtor de Picard enunciados por A.
Grothendieck en los comentarios (pp. C 07-C 11) de sus
«~Fondements de la Géométrie Algébrique~» (citado como FGA). En la primera
Exposición se demuestra el primer teorema de finitud en la forma siguiente,
más precisa que la enunciada en FGA: sean \(X\) e \(Y\) dos esquemas propios
sobre una base noetheriana integral \(S\), y sea
\(f:X\longrightarrow Y\) un \(S\)-morfismo sobreyectivo. Entonces existe un
abierto no vacío \(V\) de \(S\) sobre el cual el morfismo inducido
\(f^*:\Pic_{Y/S}\longrightarrow \Pic_{X/S}\) es representable por un
morfismo cuasiafín de presentación finita. En la segunda Exposición se
demuestra el resto.

Las dos Exposiciones difieren notablemente por sus técnicas de demostración.
En la primera se emplean de manera esencial el procedimiento de
descomposición de Oort, el descenso no plano y la representación de los
funtores no ramificados. En la segunda se reduce primero al caso en que los
esquemas \(X/S\) que intervienen son proyectivos y planos, con fibras
geométricamente integrales, o incluso, si es necesario, normales o lisos,
aplicando siempre el primer teorema de finitud a un morfismo
\(X'\to X\) adecuado. Después se emplea la herramienta de los
\((b)\)-haces creada por I. Castelnuovo y perfeccionada y utilizada por T.
Matsusaka, D. Mumford, Y. Nakai y S. Kleiman. Esta herramienta permite, en
particular, prescindir de las coordenadas de Chow, del morfismo de Frobenius
y de los teoremas de Lefschetz utilizados en las demostraciones originales
de Grothendieck.

La presente Exposición XII fue presentada oralmente, siguiendo a Grothendieck,
por M. Raynaud, a quien el redactor desea agradecer que le comunicara sus
notas personales.

\section*{1. Enunciado del teorema principal y aplicaciones}

\textbf{Teorema 1.1.} Sean \(S\) un esquema noetheriano, \(X\) e \(Y\) dos
\(S\)-esquemas propios y \(f:X\to Y\) un morfismo. Se supone que \(S\) es
integral y que \(f\) es sobreyectivo. Entonces existe un abierto no vacío
\(V\) de \(S\) tal que el morfismo
\[
  f^*:\Pic_{Y/S}|V\longrightarrow \Pic_{X/S}|V
\]
entre los funtores de Picard (localizados para la topología f.p.p.f.) sea
representable por un morfismo cuasiafín de presentación finita.

\textbf{Corolario 1.2.} Sea \(X\) un esquema propio sobre un esquema
noetheriano \(S\). Se supone que \(S\) es integral. Entonces existe un abierto
no vacío \(V\) de \(S\) tal que el funtor de Picard \(\Pic_{X/S}|_V\) sea
representable por un \(V\)-esquema en grupos que sea suma de abiertos
cuasiproyectivos de tipo finito y, por tanto, a fortiori separado sobre
\(V\).

En efecto, recordemos el hecho siguiente (FGA N\textsuperscript{o}232): sean
\(S\) un esquema y \(X\) un \(S\)-esquema proyectivo, de presentación finita
e integral sobre \(S\) (es decir, plano con fibras geométricamente
integrales). Entonces el funtor de Picard está representado por un
\(S\)-esquema en grupos \(\Pic_{X/S}\), unión creciente de abiertos
cuasiproyectivos de presentación finita. La consideración de los
«~polinomios de Hilbert~» permite aquí descomponer \(\Pic_{X/S}\) en una
suma de abiertos \(\Pic^q_{X/S}\), \(q\in\mathbb Q[t]\). En la Exposición
siguiente se demuestra, bajo la hipótesis adicional de que \(S\) sea
noetheriano, que los \(\Pic^q_{X/S}\) son de tipo finito sobre \(S\). Se sabe
entonces que 1.2 es cierto en el caso en que \(X\) es además proyectivo e
integral sobre \(S\); deduciremos de ello el caso general utilizando el lema
siguiente.

\textbf{Lema 1.3.} Sea \(X\) de tipo finito sobre \(S\) noetheriano e
integral. Entonces existen un abierto no vacío \(V\) de \(S\), un esquema
integral \(S'\) y un morfismo finito y plano \(S'\to V\), tales que toda
componente irreducible \(X'_i\) de \(X_{S'}\), provista de la estructura
reducida inducida, sea integral sobre \(S'\).

En efecto, existe una extensión finita \(L\) del cuerpo de funciones
\(k(S)\) tal que toda componente irreducible de
\((X_L)_{\text{réd}}\) sea geométricamente integral (por EGA IV 4.6.8). Sea
\(S_1\) un \(S\)-esquema integral tal que \(k(S_1)=L\) y que existan un
abierto no vacío \(V_1\) de \(S\) y un morfismo finito plano
\(h:S_1\to V_1\). Existe un abierto no vacío \(S_2\) de \(S_1\) tal que
toda componente irreducible de \(\sisi{X_{S_1}}{X_{S_2}}\), provista de la estructura % Canon SGA-SRC-001353
reducida inducida, sea integral sobre \(\sisi{S_1}{S_2}\) (EGA IV 9.7.7). Tomemos
\(V=V_1-h(S_1-S_2)\) y \(S'=h^{-1}(V)\).

Demostremos ahora 1.2. Sean \(h:S'\to V_1\) y \(X'_i\) como en 1.3. Para
todo \(i\), sea \(f_i:Z_i\to X'_i\) una presentación de Chow con \(Z_i\)
integral y \(f_i\) birracional. Tras restringir \(V\) y \(S'\) si es
necesario, se puede suponer que \(Z_i\) es integral sobre \(\sisi{S}{S'}\) (EGA IV % Canon SGA-SRC-001354
9.7.7). Sea \(f:Z=\coprod_i Z_i\longrightarrow X'=X_{S'}\) el morfismo
canónico. Entonces \(\Pic_{Z/S'}=\prod_i\Pic_{Z_i/S'}\) es representable por
un esquema en grupos que es suma de abiertos cuasiproyectivos de tipo finito.
Por último, 1.1 dice que, tras restringir \(V\) y \(S'\) si es necesario, el
morfismo \(f^*:\Pic_{X'/S'}\longrightarrow \Pic_{\sisi{Z'}{Z}/S'}\) es representable % Canon SGA-SRC-001355
por un morfismo cuasiafín. Por consiguiente, \(\Pic_{X'/S'}\) es
representable por un esquema en grupos que es suma de abiertos
cuasiproyectivos de tipo finito; gracias al descenso plano finito (SGA 1 VIII
7.6 y EGA II 6.6.1), \(\Pic_{X/S}|_V\) también lo es.

\textbf{Lema 1.4.} Sean \(k\) un cuerpo, \(P\) y \(Q\) dos
\(k\)-esquemas en grupos localmente de tipo finito, y sea \(f:P\to Q\) un
morfismo cuasiafín. Entonces \(f\) es afín.

En efecto, sea \(R=\Ker(f)\); \(R\) es un \(k\)-esquema en grupos de tipo
finito que es cuasiafín; en virtud de (SGA 3
\(\mathrm{VI}_{\mathrm B}\) 11.11), \(R\) es entonces afín. Por tanto,
\(P\to P/R\), que hace de \(P\) un \(R\)-torsor sobre \(P/R\), es afín. Por
otra parte, se sabe (SGA 3 1.4.2) que \(P/R\to Q\) es una inmersión cerrada,
y por tanto afín; en consecuencia, lo mismo ocurre con el compuesto
\(P\to P/R\to Q\).

\textbf{Corolario 1.5.} Sea \(k\) un cuerpo.
\begin{enumerate}[label=\alph*)]
\item Sea \(X\) un \(k\)-esquema propio. Entonces \(\Pic_{X/k}\) es
representable (J. P. Murre).
\item Sea \(f:X\to Y\) un morfismo sobreyectivo entre dos \(k\)-esquemas
propios. Entonces el morfismo inducido
\(f^*:\Pic_{Y/k}\longrightarrow \Pic_{X/k}\) es afín.
\end{enumerate}

\textbf{Corolario 1.6.} Sean \(S\) un esquema noetheriano integral y
\(f:X\to Y\) un \(S\)-morfismo sobreyectivo entre dos \(S\)-esquemas
propios.
\begin{enumerate}[label=\alph*)]
\item Si \(Y\) es normal sobre \(S\) (es decir, plano con fibras
geométricamente normales), entonces, tras restringir \(S\) si es necesario,
\(f^*:\Pic_{Y/S}\to\Pic_{X/S}\) es cuasifinito.
\item Si \(\underline{\Pic}_{Y/S}\) es esencialmente propio sobre \(S\) (es decir, si
satisface el criterio valuativo de ser propio, por ejemplo si \(Y\) es liso
sobre \(S\)), entonces, tras restringir \(S\) si es necesario,
\(f^*:\Pic_{Y/S}\to\Pic_{X/S}\) es finito (y, en particular, afín).
\end{enumerate}
En efecto, en a), para todo \(s\in S\), \(R_s=\Ker(f_s^*)\) es un grupo
algebraico afín; su componente neutra es un grupo algebraico propio sobre
\(k\) (FGA 236, Teorema 2.1 (iii)); por tanto, \(R_s\) es finito. En b),
\(f^*\) es a la vez propio y cuasiafín, y por tanto finito.

\textbf{Observación 1.7.} No se sabe si 1.1 sigue siendo válido cuando se
reemplaza « cuasiafín » por afín, como podrían sugerir 1.5 y 1.6.
